\section{}

We thank the reviewer for the comments. Here are points to points reply.

\comment{Abstract
\begin{itemize}
\item Overclaims novelty; FDA-based spatial modeling exists and differentiation from prior intensity-surface approaches is insufficiently clarified.
\item No explicit statement of key methodological limitations (e.g., normalization removing volume/efficiency information).
\item Practical contribution vs methodological novelty not clearly separated.
\end{itemize}
}

\reply{
  \begin{itemize}
    \item Novelty 
    \item Limitations 
    \item Contribution
  \end{itemize}

}

\comment{Introduction and Background
\begin{itemize}
	\item Motivation for density representation vs intensity-based models is weakly justified; discretization critique is not convincing since FDA also requires discretization.
	\item Limited positioning within modern representation learning or embedding-based player typologies.
	\item Missing discussion of role-dependent contextual factors (pace, lineup context, possessions, play types).
	\item Additional reference that could be added to the manuscript is: https://www.mdpi.com/2504-4990/6/3/102 and https://www.mdpi.com/2078-2489/16/8/699
\end{itemize}
}

\reply{
  \begin{itemize}
    \item Discretization 
    \item positioning
    \item Contextual factors 
    \item \textcolor{blue}{We don't see the connection between the proposed articles and ours, except they are talking about basketball. Also, we won't cite predatory journals.}
  \end{itemize}
}

\comment{
Methodology
\begin{itemize}
	\item Kernel density estimation choices (Gaussian kernel, Silverman bandwidth) insufficiently validated or sensitivity-tested.
	\item Grid resolution arbitrary without robustness analysis.
	\item Normalization removes shot volume and efficiency, limiting interpretability of components.
	\item MFPCA assumptions may violate density constraints (non-negativity, unit integral) — acknowledged but not resolved.
	\item Choice of $k=5$ clusters driven by court positions rather than data-driven criteria.
	\item Limited justification for $k$-medoids vs alternative clustering or manifold methods.
\end{itemize}
}

\reply{
  \begin{itemize}
    \item 
  \end{itemize}
}

\comment{
Results
\begin{itemize}
	\item Interpretation largely descriptive; lacks statistical validation or hypothesis testing.
	\item Low cluster agreement (ARI = 0.04–0.05) not deeply analyzed — could indicate model inadequacy rather than novel insight.
	\item Silhouette scores suggest weak cluster structure but implications are under-discussed.
	\item No external validation (e.g., predictive performance, downstream task).
\end{itemize}
}

\reply{

}

\comment{
Discussion
\begin{itemize}
	\item Claims of interpretability and practical relevance not empirically demonstrated (e.g., coaching decisions, scouting outcomes).
	\item Limited comparison against baseline models or alternative representations.
	\item Absence of multimodal or contextual variables despite acknowledging importance.
\end{itemize}
}

\reply{

}

\comment{
Conclusion and future work
\begin{itemize}
	\item Future work suggestions are broad but not technically grounded (e.g., multimodal integration lacks methodological direction).
	\item Does not address scalability, computational complexity, or real-time applicability.
	\item Key methodological limitation (density normalization) insufficiently emphasized.
\end{itemize}
}

\reply{

}

\comment{
Technical Writing
\begin{itemize}
	\item Some redundancy in explaining FPCA concepts.
	\item Overly descriptive narrative in results; could benefit from more concise statistical framing.
	\item Terminology occasionally inconsistent between “shooting frequency,” “density,” and “intensity.”
	\item Figures interpretation sometimes qualitative rather than quantitative.
	\item Keywords are not precise and should be in alphabetical order.
\end{itemize}
}

\reply{
  \begin{itemize}
    \item Some repetition are intentional because the paper sits at the intersection of functional data analysis and sports analytics. Not all readers will be equally comfortable with FDA terminology, so moderate repetition improves accessibility.
    \item Desscriptive results
    \item Terminology
    \item Figures
    \item The keywords are in alphabetical order. The revised keywords are : Basketball analytics, Clustering, Functional data analyasis, Shot charts and Shot-density estimation.
  \end{itemize}
}


